一份上线评审里写着：``该接口的平均延迟是 120 毫秒。''

数字没错，算法也没错。上线之后投诉照来——因为真正让用户放弃的不是平均那一次请求，是最慢的那百分之一。而这个分布的尾巴很重，重到平均值几乎不携带关于尾部的信息：把 P99 从 900 毫秒压到 400 毫秒，平均值可能只动几毫秒。

这类事故的根源不在计算。它在于交出结论的人没说清自己算的是什么、这个数字在什么条件下能支撑什么判断。概率论真正难的部分正是这一段：随机实验到底是什么、哪些独立性假设来自现实而哪些只是为了好算、以及手上这个结论是精确概率、保守上界，还是一个只在样本足够多时才成立的近似。

最后这个区分尤其容易被忽略，所以先摆出来。

\begin{center}
\begin{tikzpicture}[x=1cm,y=1cm,font=\small,>=stealth]
  \draw[black!45,fill=green!8] (0,2.2) rectangle (13.2,3.5);
  \node[align=left,anchor=west] at (0.25,3.1) {\textbf{精确分布}\quad 知道生成机制，概率可以直接算出来};
  \node[align=left,anchor=west,font=\footnotesize,black!70] at (0.25,2.55) {代价：模型必须真的对。机制写错时，算得再精确也只是精确地错};
  \draw[black!45,fill=blue!8] (0,0.6) rectangle (13.2,1.9);
  \node[align=left,anchor=west] at (0.25,1.5) {\textbf{有限样本界}\quad Markov、Chebyshev、Hoeffding};
  \node[align=left,anchor=west,font=\footnotesize,black!70] at (0.25,0.95) {对任意样本量都成立，且只需很弱的假设。代价：往往保守得多};
  \draw[black!45,fill=orange!10] (0,-1.0) rectangle (13.2,0.3);
  \node[align=left,anchor=west] at (0.25,-0.1) {\textbf{渐近近似}\quad 大数定律、中心极限定理};
  \node[align=left,anchor=west,font=\footnotesize,black!70] at (0.25,-0.65) {\(n\) 大时很准，但不告诉你多大才够；而且尾部最不准，偏偏尾部最要紧};
  \draw[->,black!50] (13.5,3.0) -- (13.5,-0.5);
  \node[rotate=-90,black!65,font=\footnotesize] at (13.95,1.25) {假设越来越弱，结论越来越松};
\end{tikzpicture}

\emph{三种结论的强度完全不同。实践中最常见的事故，是把第三种当成第一种用——拿一个渐近近似去回答尾部问题。}
\end{center}

\subsection{五条贯穿始终的判断}

\textbf{多数概率争论不是算错，是两个人心里的样本空间不一样。}三门问题、检验悖论、幸存者偏差，还有日常的``日活到底按用户算还是按会话算''，追到底都是 \(\Omega\) 没说清。概率的建模三件套 \((\Omega,\mathcal F,P)\) 看着形式化，其实每一件都在逼你交代一件容易含混的事：结果有哪些、你能观测到哪些、以及怎样给它们指派数。

\textbf{条件概率是把样本空间截到已知信息之内再重新归一。}由此立刻可知 \(P(A\given B)\) 和 \(P(B\given A)\) 是两回事，混淆二者就是检察官谬误。Bayes 公式的价值不在代数形式，而在它说明了后验由先验和似然比共同决定：患病率 \(0.1\%\) 时，一个灵敏度和特异度都是 \(99\%\) 的检验，阳性结果的实际患病概率还不到 \(10\%\)。罕见事件检测里的假阳性淹没现象，全部来自这一条。

\textbf{独立与条件独立互不蕴含，两个方向都不成立。}共同原因会让两个本无关系的量相关，控制住原因后关联消失；共同结果反过来——两个独立的量，一旦条件在它们共同导致的结果上就会产生相关。后者是选择偏差的数学形态，也解释了为什么控制变量不是越多越好。

\textbf{分布由机制生成，不该当成表来背。}Poisson、指数、Gamma 不是三个分布，是同一个泊松过程的三种切法；几何分布与二项分布是同一套试验换个提问方向。认出机制就能预判性质，背下公式则一换场景就失效。更要紧的是知道机制什么时候不成立：重尾分布下均值甚至方差可以根本不存在，那时``平均值会稳定下来''这句话是错的。

\textbf{极限定理是有前提的，前提塌了结论就塌。}大数定律要求期望存在，中心极限定理还要求方差有限。Cauchy 分布的样本均值不收敛——不是收敛得慢，是压根不收敛。开头那个延迟的例子正在这条边界上：重尾数据里，样本均值本身就是个不稳定的统计量。

\subsection{八章怎么安排}

\hyperref[ch:037]{``概率空间：先说清楚什么可能发生''}和\hyperref[ch:038]{``条件概率、Bayes 与独立性：新信息怎样改变判断''}是地基，建议连读。它们几乎不含计算，但决定了后面所有模型写得对不对。

\hyperref[ch:039]{``随机变量与分布：把复杂结果压缩成可研究的数值''}换一种记账方式：把结果映到数轴上，从此可以整体地谈一个分布，而不必一个事件一个事件地算。\hyperref[ch:040]{``期望与矩：怎样用少量数字概括一个分布''}接着讨论要保留哪些摘要——其中``条件期望是最佳预测''一篇价值最高，它就是回归的目标，也是投影在概率里的化身。

\hyperref[ch:041]{``联合分布：单个分布之外，变量怎样一起变化''}处理多个量的共同变化；边缘分布相同不足以确定依赖结构，协方差也只看得见线性关系。\hyperref[ch:042]{``常见分布：不要背表，要先辨认随机机制''}适合当作手册按需查阅，但其中重尾与极值两篇请完整读一遍，它们改变的是判断而不是计算。

\hyperref[ch:043]{``变换与概率界：不知道完整分布时还能说什么''}和\hyperref[ch:044]{``极限定理：重复随机性为何会产生稳定规律''}收尾。前者给出一条精度阶梯：交出的信息越多，界越紧；后者说明重复随机为何产生稳定，以及这种稳定在什么条件下不再出现。做实验设计、估样本量、判断置信区间可不可信，用的都是这两章。

读的时候有个动作值得保持：每写下一个概率乘积，问一句它的来历——是定义带来的恒等式，是实验设计（随机化）保证的，还是仅仅因为这样写最省事。前两种可以依靠，第三种是需要标注的假设，而它恰恰是最常见的那种。
